\chapter{Geometry at One Base Point}
\label{comp:ch02}

\section{The question and the answer}
\label{comp:ch02-question}

The technical chapter asks what geometric data an agent carries at the single base point $c_*$ at which the
whole main manuscript works, and which of those data are invariant, meaning unchanged by any change in how
the data are written down. It gives three answers. The group of genuine symmetries at $c_*$, the active
gauge group, is one copy of the structure group $G$ sitting inside the much larger group $\prod_{i\in V}G$
of independent per-agent frame changes, which the chapter calls passive (\MT{prop:geo-active-image-point}).
Every independent link between two agents carries a conjugacy class that neither passive nor active changes
can alter, and the holonomy of a closed loop of agents is the telescoped product of the per-edge twists whose
classes these are (\MT{prop:geo-edge-twist-invariance}); the classes alone do not determine the holonomy.
Three notions of agreement on a block of agents, frame agreement, agreement of an equivariant statistic
such as the mean, and law coherence, are logically distinct, and the chapter proves exactly which
implications hold among them (\MT{prop:geo-three-agreements}). These objects are used downstream: passive
re-expression is the transformation under which the generative law must be covariant (\cref{comp:ch04}),
link-transported comparisons are what attention compares (\cref{comp:ch06}), and holonomy together with the
trivialization criterion becomes the fixed sectors of coarse channels (\cref{comp:ch10}).

Is the principal- and associated-bundle machinery doing real work at one base point? Only as typing. At
$c_*$ the principal bundle enters only through its fiber $(\mathscr P_G)_{c_*}$, which is a $G$-torsor,
that is, a copy of $G$ with no preferred identity element; the belief and model bundles enter only through
their fibers, which are $G$-spaces of probability laws; and a section enters only through its value at
$c_*$, which is a law. Nothing in the chapter uses a connection, a base curve, the topology of
$\mathscr P_G$, or the manifold structure of the context space $\mathcal C$, and the manuscript says so
itself: graph flatness is a property of the links alone, whatever the topology of $\mathscr P_G$
(\MT{rem:geo-flatness-strictness}), and graph-link holonomy is not the holonomy of any base connection
(\MT{sec:intro-conventions}). What the bundle language contributes is a clean separation between coordinates
(frames, link coordinates, sample coordinates) and frame-independent objects (link maps, points of the
associated fibers), and both unfamiliar statements of the chapter come out of that separation. The link sector
itself is lattice gauge theory on the interaction multigraph, with the lattice vertex law, Wilson loops, and the
lattice criterion for pure gauge; what is added is a group-valued frame at every agent, all frames lying in one
shared fiber. Transport along the base, which would make the bundle language indispensable, is deferred to
Supplement B.

\section{The picture}
\label{comp:ch02-picture}

\subsection{One fiber, many frames}

All agents sit at the same base point, so their frames $u_i$ are points of one and the same fiber. Pick any
reference point $\sigma$ of that fiber and each frame becomes a group element $U_i$ through
$u_i=\sigma\cdot U_i^{-1}$; any two frames are then compared by $T_{ij}=U_iU_j^{-1}$, which does not depend
on $\sigma$ (\MT{prop:geo-frame-comparisons}). A useful picture is a set of observers in one laboratory,
each with private coordinate axes. An observer who changes axes, a passive re-expression $g_i$, changes
every description: sample coordinates are multiplied by $R_i^x=\rho_x(g_i)$, laws are pushed forward, the
comparisons become $g_iT_{ij}g_j^{-1}$, and nothing in the laboratory has moved. Rotating the laboratory
itself, an active transformation, that is, an automorphism of the fiber that commutes with the right action
of $G$, is a single operation, which observer $i$ records as $\varphi_i$ and observer $j$ as the conjugate
$\varphi_j=T_{ij}^{-1}\varphi_iT_{ij}$. The active transformations therefore form a copy of $G$, of
dimension $\dim G$ rather than $|V|\dim G$, and a tuple that changes the frame of one agent alone is not
active when $G$ is nontrivial and there are at least two agents (\MT{prop:geo-active-image-point}).

\begin{analogy}{structural}
In lattice gauge theory a gauge transformation has two readings that give the same group: a change of local
trivialization at every site, and a vertical automorphism of the bundle over the lattice. Both are
$\prod_xG$, because every site $x$ has its own fiber and nothing but the link variable relates the fibers at
two sites; Elitzur's theorem \citep{elitzur1975impossibility} concerns this local group. At $c_*$ the two
readings come apart. The passive group of frame changes is still $\prod_{i\in V}G$, and it acts on link
coordinates by exactly the lattice formula $\Theta_e\mapsto g_i\Theta_eg_j^{-1}$
(\MTeq{eq:geo-regime-two-gauge-law}). The automorphisms form only one copy of $G$, because the agents are not
sites with fibers of their own: they are frames in one shared fiber, and every automorphism of the whole
bundle restricts at $c_*$ to one element of that copy (\MT{prop:geo-active-image-point}). Lattice sites
correspond to base points, not to agents at one base point; which reading a lattice model of the supplements
has is left open there (\SB{open:rg-lattice-typing}). The correspondence is structural: the same
transformation formula for the links, a different automorphism group.
\end{analogy}

\subsection{A link is an automorphism of the shared fiber}

A link $\Theta_e^x\in G$ on the oriented edge copy $e:j\to i$ defines the link map
$\mathsf t_e^x(u_j\cdot g)=u_i\cdot\Theta_e^xg$, a right-equivariant bijection of the one fiber onto itself
that does not depend on the frames used to write it (\MT{prop:geo-edge-twist-invariance}). An equivariant
self-map of a torsor is, once a reference point is chosen, left multiplication by one group element, here
the per-edge twist $\mathring\Theta_e^x=U_i^{-1}\Theta_e^xU_j$. Passive re-expressions leave every twist
unchanged, a change of reference point conjugates all twists by one common element, and an active
transformation does the same. Each edge copy therefore carries the invariant conjugacy class
$[\mathring\Theta_e^x]$, and the full invariant content of the links is the twist tuple modulo one
simultaneous conjugation. In lattice gauge theory a link is a map between two different fibers and has no
conjugacy class; here all agents share a fiber, so a link is a symmetry of that fiber and has one.

In lattice language, and informally, the twist is the link measured against the direct frame comparison.
The comparison $T_{ij}$ transforms under passive re-expression exactly like a link from $j$ to $i$, so
$T_{ij}^{-1}\Theta_e^x$ is the holonomy of a two-step loop based at $j$: out to $i$ along the link, back
along the frame comparison. It is conjugate to the twist. The per-edge class is thus a Wilson loop, of a
loop that closes through the frame comparison instead of through another link, and pure lattice gauge
theory has no such loops because it has no frame comparison between sites. The same reading removes an
apparent clash with lattice intuition. A single link coordinate between two distinct agents carries no
invariant by itself, since the re-expression $g_i=e_G$, $g_j=\Theta_e^x$ sends it to $e_G$
(\MT{rem:geo-flatness-strictness}), exactly as in lattice gauge theory. The invariant belongs to the pair
formed by the link and the frames.

Whether a model can see these extra invariants depends on whether it reads the frames. Informally: if a model
satisfies the kernel-covariance hypothesis of \cref{comp:ch04} with parameters that the re-expression leaves
unchanged, and none of its kernels reads the frames, then two link configurations related by the lattice
vertex law, with the frames held fixed, have the same evidence, by the covariance of the joint
(\MT{prop:gen-product-evidence-invariance}). Such a model sees only the Wilson-loop data, and the running
example below is of this kind. A mechanism that contains a map shared by all agents and referred to their
frames reads the twists directly; the shared query and key reference maps of attention are of this kind
(\MT{tab:obs-transformer-correspondence}). The per-edge classes are physical to the extent that the model
contains frame-referenced structure.

\begin{analogy}{structural}
Write the frames as group elements $U_i$ against one reference point. A passive re-expression acts by
$U_i\mapsto g_iU_i$, a change of reference point by $U_i\mapsto U_ig_0$, and the twist $U_i^{-1}\Theta_eU_j$
is invariant under the first and conjugated by the second. These are the field content and the symmetry
pattern of a lattice gauge--Higgs model whose Higgs field is group valued (for $SU(2)$, a fundamental doublet
of fixed length written as a matrix): the twist is the link in unitary gauge, and the single remaining global
group acts on unitary-gauge links by conjugation \citep{FradkinShenker1979,OsterwalderSeiler1978}.
Supplement B spells out this reading for $G=\mathbb Z_2$, identifying the matter sign with the frame
coordinate and the dressed link sign with the twist (\SB{open:rg-lattice-typing}). The correspondence is
structural and kinematic only. The frames are structural data at $c_*$, not a field that is summed over, so
nothing about gauge--Higgs dynamics, phase structure, or complementarity carries over.
\end{analogy}

\subsection{Pure gauge, frame alignment, and Wilson loops}

Links are graph flat on a block $I$ when every closed internal walk has trivial product, equivalently when they
are a coboundary $\Theta_e^x=V_iV_j^{-1}$ for vertex potentials $V_i$, unique up to one right multiplication per
connected component (\MT{prop:geo-trivializing-criterion}, \MT{def:geo-graph-flat-links}). This is pure gauge.
Links are frame aligned when $\Theta_e^x=T_{ij}$, the canonical pure-gauge configuration supplied by the frames
themselves; then every twist is $e_G$ and every link map is the identity of the fiber. Frame alignment implies
graph flatness and not conversely, and on a tree every configuration is graph flat while its per-edge classes can
all be nontrivial (\MT{rem:geo-flatness-strictness}). In lattice terms: on a tree all links are pure gauge, as in
the maximal-tree gauge fixing \citep{Creutz1977}, and here they still carry twists relative to the frames.

\subsection{Three kinds of agreement}

On a connected block $I$ the chapter separates three conditions (\MT{def:geo-law-coherence}). Frame
agreement is graph flatness of the links on $I$, a condition on the links alone. Agreement of an equivariant
statistic, with the mean as the basic case, asks that transported statistics match on every internal edge,
$\mu_i=\rho_x(\Theta_e^x)\mu_j$, a condition on links and means. Law coherence asks that whole laws match,
$q_i=\widehat\rho_x(\Theta_e^x)q_j$, a condition on links and laws. The relations are proved in
\MT{prop:geo-three-agreements}. Law coherence holds exactly when the base law $q_{i_0}$ is fixed by every
loop holonomy based at $i_0$ and the other laws are its transports, and mean agreement holds exactly when
the base mean lies in the common fixed space of the represented holonomies and the other means are its
transports. Law coherence implies agreement
of every equivariant statistic, and the converse fails for the mean. Frame agreement neither implies nor is
implied by either of the other two. With several independent cycles, a condition checked on one loop is
necessary and not sufficient.

The distinctions matter because each notion is a candidate criterion for when agents form a collective,
and the three give different verdicts. Frame agreement without law coherence is a block with no flux whose
agents simply disagree. Law coherence without frame agreement is a block with nontrivial flux whose laws
agree anyway, because the flux lies in the stabilizer of the common law: isotropic Gaussians around the
running example's loop, or, for $G=\GL(2,\R)$, the law $\mathcal N(0,\Sigma)$ transported around a loop
whose holonomy $\Sigma^{1/2}R(0.7)\Sigma^{-1/2}$ is not orthogonal and fixes no nonzero vector. Mean
agreement without law coherence is a block whose covariances disagree. The coarse-graining chapters use the
third notion: they score a block by its divergence from holonomy-fixed parent laws, and an attained zero score
means exactly that the transported member laws coincide in one such law, which makes them law coherent
(\MT{thm:cg-holonomy-kl-marginal}). A collective criterion built on frame agreement would be too strong for that
purpose, and one built on means too weak. None of the three selects a partition into meta-agents,
which is part of \MT{open:cg-natural-selector}.

\begin{analogy}{exact}
Part (i) of \MT{prop:geo-three-agreements} is the graph version of the holonomy principle of differential
geometry: a covariantly constant section through a given fiber element exists exactly when that element is
fixed by the holonomy group. Law coherence is a covariantly constant section of the law bundle over the
block, with the links playing the connection.
\end{analogy}

\begin{analogy}{structural}
A covariantly constant Higgs field forces the holonomy of the gauge field into the stabilizer of the Higgs
value, the unbroken subgroup, while flux inside that subgroup is allowed. Law coherence with nontrivial
holonomy has the same pattern with a probability law in place of the Higgs value: flux is allowed exactly
inside the law's stabilizer. No symmetry breaking is claimed by the chapter.
\end{analogy}

\subsection{Compact and noncompact structure groups}

Nearly everything in the chapter holds for an arbitrary Lie group $G$, and the trivialization criterion
holds for an arbitrary group, so the frame comparisons, the single active copy of $G$, the per-edge classes,
and the three agreements hold verbatim for $\GL(K,\R)$, $\GL(K,\mathbb C)$, and $SU(N)$. The exception is
the construction of a probability law on the links from loop penalties (\MT{rem:geo-loop-penalties}). For a
unitary representation the penalty $1-\operatorname{Re}\Tr\mathsf U/n$ of a represented loop holonomy
$\mathsf U$ is Wilson's plaquette action \citep{WilsonConfinement1974} normalized per dimension; it lies in
$[0,2]$, and for compact $G$ the Boltzmann weight of finitely many such penalties against the product Haar
probability measure on the links has a normalizer between $e^{-2\sum_p\mathrm c_p}$ and $1$. This covers
$SU(N)$ and $SO(K)$ with unitary representations. For $\GL(n,\R)$ with $n\geq2$ the same remark proves a
no-go: no norm on $\mathfrak{gl}(n,\R)$ is invariant under conjugation by all of $\GL(n,\R)$. An invariant
nonnegative penalty vanishing exactly at the identity exists once a symmetric positive-definite fiber metric
at the base vertex is declared and made to transform covariantly, which is additional data, playing the role
a metric plays in reducing a general linear group to an orthogonal one. Whether any such penalty gives a
normalizable law on $\GL$-valued links is declined in the chapter's scope paragraph.

The gap is not a technicality. Informally, the Haar measure of a noncompact group is infinite, and a penalty
that depends only on loop holonomies leaves every tree direction of the link configuration unpenalized, so
the naive Boltzmann weight already fails to integrate on a graph with a single edge. A proper reference that
is not Haar would normalize such a law at the cost of the invariance the penalty was built to have; the
manuscript makes that trade once, for the Gaussian presentations of its formation model, as recorded in
\MT{open:prob-law-point-prior}. Two further facts bear on the target group. Links are structural parameters
by default, not random variables (\MT{sec:intro-conventions}), so the missing $\GL$ link law blocks only models
that integrate over links. And for nondegenerate Gaussian laws under independent $\GL(K)$ links, law coherence
around a loop requires the holonomy to be conjugate into $\mathrm O(K)$, which forces $\lvert\det\rvert=1$ on
the loop and so holds only on a Lebesgue-null set of link configurations
(\MT{cor:cg-holonomy-noncompact-gaussian}, discussed in \cref{comp:ch10}); for the target group, law coherence of
Gaussian beliefs around a loop is nongeneric.

\section{Objects and notation}
\label{comp:ch02-objects}

\begin{decoder}
Base point $c_*$; context space $\mathcal C$ & The single point of the context manifold at which everything happens; $\mathcal C$ itself is idle in this chapter. & One spacetime point, or one lattice site \\
Principal bundle $\mathscr P_G$; its fiber at $c_*$ & At $c_*$ only the fiber matters: a $G$-torsor, a copy of $G$ without a preferred identity. & Space of frames at one point \\
Frame $u_i$; reference point $\sigma$; $U_i$ & Agent $i$'s chosen point of the shared fiber; against a reference $\sigma$ it is the group element $U_i$ with $u_i=\sigma\cdot U_i^{-1}$. & Local frame; group-valued Higgs field in the analogy above \\
Passive re-expression $g_i$; passive vertex group $\prod_{i\in V}G$ & Agent $i$ replaces its frame by $u_i\cdot g_i^{-1}$; coordinates change, nothing else does. & Local gauge transformation read as a change of trivialization \\
Active gauge group $\mathcal A_{c_*}$; local functions $\varphi_i$ & Automorphisms of the shared fiber; one copy of $G$, written in agent $i$'s frame as $\varphi_i$. & One global symmetry transformation \\
Frame comparison $T_{ij}$ & The unique group element with $u_j=u_i\cdot T_{ij}$, equal to $U_iU_j^{-1}$. & Canonical pure-gauge link on the complete graph \\
Sample fibers $\mathsf K$, $\mathsf M$; representations $\rho_b$, $\rho_m$ & Spaces of belief samples and model presentations, with $G$ acting; the two actions may be inequivalent, unfaithful, and of different dimension. & Matter representations \\
Law fibers $\mathcal B_b$, $\mathcal B_m$; pushforward action $\widehat\rho_x$ & Families of probability laws on the sample fibers, on which $G$ acts by pushforward. & A nonlinear target space whose points are laws \\
Associated bundles $\mathcal E_b$, $\mathcal E_m$ & Bundles whose fiber at $c_*$ is a space of laws, identified with $\mathcal B_x$ by each frame. & Matter bundle, with laws in place of field values \\
Agent $\mathcal A^i=(\mathcal C_i;q_i^b,q_i^m,u_i)$ & A frame plus a belief law and a model law; at $c_*$ only their values matter. & None; a naming decision \\
Cross morphisms $\Phi$, $\widetilde\Phi$ & Fiberwise maps from belief laws to model laws and back; extra structure, not supplied by the bundle. & None \\
Multigraph $\mathfrak I=(V,E,\bar{\ })$; edge copy $e:j\to i$; reversal $\bar e$ & Oriented edges between agents, parallel copies allowed; every copy has a reversed partner. & Lattice with multiple bonds \\
Link $\Theta_e^x$; vertex law & A group element per edge copy and channel, independent of the frames; transforms by $g_i\Theta_e^xg_j^{-1}$. & Link variable $U_{xy}$ \\
Link map $\mathsf t_e^x$ & The frame-independent automorphism of the fiber that the link defines. & None in pure lattice gauge theory \\
Per-edge twist $\mathring\Theta_e^x$ & $U_i^{-1}\Theta_e^xU_j$, the link written against the single reference point; its conjugacy class is invariant. & Link in unitary gauge \\
Path product $\Theta_p^x$; holonomy $\mathrm{Hol}^x(p)$ & Ordered product of links along a walk; for a closed walk, a group element covariant under conjugation. & Wilson loop, before taking the class \\
Graph flat; jointly trivializing & Trivial holonomy on every internal loop; in both channels at once. & Pure gauge; zero field strength \\
Frame aligned & Links equal the frame comparisons, so all twists are trivial. & The canonical pure-gauge background \\
Law coherent & Laws on a block equal each other's transports along every internal edge. & Covariantly constant matter field \\
$s$-agreeing; mean agreement & An equivariant statistic, such as the mean, agrees after transport. & Agreement of an order parameter \\
\end{decoder}

Several words and glyphs trip a physicist. "Gauge" appears with two qualifiers: the passive group is the local
one and the active group is global, and "gauge invariant" without qualifier means invariant under both.
"Holonomy", "flat", and "twist" are graph-link notions: a flat link configuration says nothing about the
curvature of any connection on $\mathcal C$, and a twist is not a twist of the bundle, which over one point is
always trivial. "Common frame" means one frame shared by an agent's belief and model channels
(\MT{prop:geo-common-frame-gauge}), not one frame shared by all agents. The orientation runs from source to
receiver: $e:j\to i$ carries $j$'s coordinates into $i$'s, a path product is written so that its rightmost link
acts first, and $T_{ij}$ transforms like a link on $e:j\to i$. $T_{ij}$ is never a temperature or a transfer
matrix, a hat on $\rho$ means pushforward, and a sans-serif $\mathsf G$ is a metric, not the group. Calling a
section-bearing object an agent is a naming decision with no claim of agency attached
(\MT{sec:intro-conventions}).

\section{The formulas, explained}
\label{comp:ch02-formulas}

\paragraph{Frames and comparisons.}
The fiber at $c_*$ is a right $G$-torsor, and the chapter's first relations are
\[
u_i=\sigma\cdot U_i^{-1},
\qquad
u_j=u_i\cdot T_{ij},
\qquad
T_{ij}=U_iU_j^{-1},
\qquad
T_{ij}'=g_iT_{ij}g_j^{-1},
\qquad
U_i'=g_iU_i
\]
(\MT{prop:geo-frame-comparisons}). The reference point $\sigma$ is an arbitrary origin, and $T_{ij}$ does not
depend on it. The identities $T_{ii}=e_G$, $T_{ji}=T_{ij}^{-1}$, $T_{ij}T_{jl}=T_{il}$ make $(T_{ij})$ a pure-gauge
configuration on the complete graph. The re-expression $g_i=U_i^{-1}$ gives every agent the frame $\sigma$ and
all $T_{ij}=e_G$; this equal-frame gauge is always reachable at $c_*$, and the running example uses it.

\paragraph{The active image.}
An automorphism $F$ of the fiber, written in agent $i$'s frame as $F(u_i)=u_i\cdot\varphi_i$, satisfies
$\varphi_j=T_{ij}^{-1}\varphi_iT_{ij}$ (\MTeq{eq:geo-active-image-point}), so the whole tuple is fixed by any one
entry. In coordinates it acts like the passive re-expression $g_i=\varphi_i$; the difference is that the active
map moves the configuration while the passive one relabels it.

\paragraph{Links, the vertex law, and holonomy.}
For each edge copy $e:j\to i$ and channel $x\in\{b,m\}$,
\[
\Theta_{\bar e}^x=(\Theta_e^x)^{-1},
\qquad
\Theta_e^{x\prime}=g_i\Theta_e^xg_j^{-1},
\qquad
\mathrm{Hol}^x(p)=\Theta_{e_0}^x\Theta_{e_1}^x\cdots\Theta_{e_{r-1}}^x,
\qquad
\mathrm{Hol}^x(p)'=g_{i_0}\mathrm{Hol}^x(p)g_{i_0}^{-1},
\]
the last two for a closed walk based at $i_0$ (\MT{def:geo-graph-links}). These are the lattice formulas term
for term. The two channels carry separate link families, and the common-link case $\Theta_e^b=\Theta_e^m$ is a
declared specialization.

\paragraph{Twists.}
With $u_i=\sigma\cdot U_i^{-1}$,
\[
\mathring\Theta_e^x=U_i^{-1}\Theta_e^xU_j,
\qquad
\mathrm{Hol}^x(p)=U_{i_0}\mathring\Theta_{e_0}^x\mathring\Theta_{e_1}^x\cdots\mathring\Theta_{e_{r-1}}^xU_{i_0}^{-1}
\]
(\MTeq{eq:geo-edge-twist}, \MTeq{eq:geo-edge-twist-holonomy}). The twist is the link rewritten against the
single reference point $\sigma$; frame factors cancel in pairs along a closed walk. For abelian $G$ each twist is
itself invariant (Supplement B's $\mathbb Z_2$ reading calls it the dressed link sign), and on a tree the twists
are the entire invariant content.

\paragraph{The trivialization criterion.}
For group elements $\Xi_e$ on the internal edge copies of a block,
\[
\begin{aligned}
&\Xi_e=V_iV_j^{-1}\ \text{for every internal }e:j\to i\\
\Longleftrightarrow\quad&\Xi_{e_0}\Xi_{e_1}\cdots\Xi_{e_{r-1}}=e_G\ \text{for every closed internal walk}
\end{aligned}
\]
(\MTeq{eq:geo-trivializing-criterion}). The first line is pure gauge and the second is triviality of all Wilson
loops. The proof uses only the group law, so it holds for $\GL(K)$, for its subgroups, and for represented
links, with no compactness or topology.

\paragraph{Loop penalties.}
For a represented loop holonomy $\mathsf U=\rho_x(\mathrm{Hol}^x(p))$ of dimension $n$,
\begin{gather*}
\frac{1}{2n}\Tr\bigl((I_n-\mathsf U)^*(I_n-\mathsf U)\bigr)=1-\frac1n\operatorname{Re}\Tr\mathsf U\in[0,2]
\qquad(\mathsf U\text{ unitary}),\\
\Tr\bigl((\mathsf G^{\mathrm{fib}})^{-1}(I_n-\mathsf U)^\top\mathsf G^{\mathrm{fib}}(I_n-\mathsf U)\bigr)
\qquad(\mathsf U\text{ real}).
\end{gather*}
The first is Wilson's action per loop, vanishing exactly at $\mathsf U=I_n$; for $U(1)$ it is $1-\cos\phi$. It
is invariant because a passive re-expression conjugates $\mathsf U$ by the unitary $\rho_x(g_{i_0})$. The second
is the general linear substitute (\MT{rem:geo-loop-penalties}); it is invariant only because the fiber metric
$\mathsf G^{\mathrm{fib}}$ is declared to transform by
$(R_{i_0}^x)^{-\top}\mathsf G^{\mathrm{fib}}(R_{i_0}^x)^{-1}$, so it is a covariant field at the base vertex,
not a constant.

\paragraph{Agreement.}
Law coherence and its mean-level shadow read
\[
q_i=\widehat\rho_x(\Theta_e^x)q_j\ \text{on each internal }e:j\to i,
\qquad
\mu_{i_0}\in\bigcap_{p\text{ closed at }i_0}\ker\bigl(\rho_x(\mathrm{Hol}^x(p))-\operatorname{id}\bigr)
\]
(\MTeq{eq:phil-law-coherence}, \MTeq{eq:phil-agreement-kernel}). The first is a parallel section of the law
bundle over the block; the second says the base mean survives every loop. For a nontrivial planar rotation loop
the kernel is $\{0\}$ and every agreeing mean is zero.

\paragraph{Cross morphisms.}
A cross morphism has representatives transforming by
$\Phi_i'=\widehat\rho_m(g_i)\circ\Phi_i\circ\widehat\rho_b(g_i)^{-1}$, and one map $\Phi_0$ serves in every frame
exactly when it intertwines the two actions (\MT{prop:geo-intertwining-cross-map}). In the running example's
setting, with $SO(3)$ acting on $\R^3$ in the belief channel and trivially on $\R$ in the model channel, the only
equivariant linear map $\R^3\to\R$ is zero by Schur's lemma, while pushforward of laws by the norm,
$\beta\mapsto(\lvert\cdot\rvert)_\#\beta$, is an equivariant map of laws: law-level equivariance is weaker than
linear equivariance.

\paragraph{The running example.}
Take $G=SO(3)$, $\rho_b$ the defining representation on $\R^3$, $\rho_m$ trivial on $\R$, three agents on a
triangle with $e_1:2\to1$, $e_2:3\to2$, $e_3:1\to3$, all frames equal, and belief links
\[
\Theta_{e_1}^b=R_z(\tfrac\pi2),
\qquad
\Theta_{e_2}^b=R_x(\tfrac\pi2),
\qquad
\Theta_{e_3}^b=I_3,
\qquad
\mathrm{Hol}^b(p)=R_z(\tfrac\pi2)R_x(\tfrac\pi2),
\qquad
\Tr\mathrm{Hol}^b(p)=0
\]
(\MT{sec:geo-running-example}, \MTeq{eq:geo-running-holonomy}). The per-edge classes are the rotation angles
$\pi/2$, $\pi/2$, $0$, and the holonomy is a rotation by $2\pi/3$, here about the axis $(1,1,1)$. Replacing the
second quarter turn by one about the same axis as the first gives a rotation by $\pi$ with the same three
classes, which is the concrete reason the class tuple alone does not determine the holonomy while the tuple
modulo one simultaneous conjugation does. The loop is not graph flat, yet the isotropic laws
$\mathcal N(0,s^2I_3)$ are law coherent on it, because every rotation fixes them. The one record of the
example has success probability $\exp(-\tfrac12\lvert k_1-\rho_b(\Theta_{e_1}^b)k_2\rvert^2)$, a transported
discrepancy whose length no passive re-expression changes (\MTeq{eq:geo-running-record}). In lattice terms,
this is a frustrated plaquette with flux $2\pi/3$ and a matter law that does not feel the flux.

\section{What is proved, and how solid it is}
\label{comp:ch02-status}

\begin{scorecard}
Agents and law bundles: \MT{def:geo-associated-bundles}, \MT{def:geo-agent} & An agent is a frame in the one fiber plus a belief law and a model law; the channels may carry inequivalent representations. & \CHOICE & Standard construction; the commitment is laws valued in $G$-spaces at one fiber. \\
Frame comparisons: \MT{prop:geo-frame-comparisons} & $T_{ij}$ is unique, a cocycle, and transforms by $g_iT_{ij}g_j^{-1}$. & \SOLID & Any Lie group; torsor algebra. \\
Active image: \MT{prop:geo-active-image-point} & Fiber automorphisms form one copy of $G$ inside $\prod_iG$, proper when $G\neq\{e_G\}$ and $|V|\geq2$. & \SOLID & Any Lie group; standard in substance. \\
Common frame: \MT{prop:geo-common-frame-gauge} & Separate belief and model frames reduce to one by a re-expression. & \SOLID & Harmless because nothing reads the relative frame. \\
Cross morphisms: \MT{prop:geo-intertwining-cross-map} & Gauge law of representatives; intertwiners are exactly the frame-independent ones; not every cross morphism is induced; Schur vanishing does not remove cross morphisms. & \SOLID & Their existence is additional structure. \\
Independent links: \MT{def:geo-graph-links} & One group element per oriented edge copy and channel, independent of the frames. & \CHOICE & Structural parameters by default. \\
Twist invariance: \MT{prop:geo-edge-twist-invariance} & Each edge carries an invariant twist class; link sets are actively related exactly when their twist tuples are simultaneously conjugate. & \SOLID & Any Lie group; standard in substance; an invariant of links with frames. \\
Trivialization: \MT{prop:geo-trivializing-criterion} & Coboundary exactly when all closed-walk products are trivial; potentials unique up to one right factor per component. & \SOLID & Any group; pure gauge exactly when Wilson loops are trivial. \\
Strictness of flatness: \MT{rem:geo-flatness-strictness} & Frame alignment is strictly stronger than graph flatness; one link coordinate has no invariant; trees can carry nontrivial classes. & \SOLID & General non-implications on minimal graphs. \\
Compact loop penalty: \MT{rem:geo-loop-penalties} & For unitary $\rho_x$, $1-\operatorname{Re}\Tr\mathsf U/n\in[0,2]$ is invariant; for compact $G$ the penalty law against product Haar probability has normalizer in $[e^{-2\sum_p\mathrm c_p},1]$. & \COND & Needs $G$ compact, $\rho_x$ unitary: holds for $SU(N)$, $SO(K)$; fails for $\GL(K,\R)$, $\GL(K,\mathbb C)$. \\
No invariant norm for $\GL$: \MT{rem:geo-loop-penalties} & No norm on $\mathfrak{gl}(n,\R)$, $n\geq2$, is $\GL(n,\R)$-conjugation invariant; a covariant fiber metric restores an invariant penalty. & \SOLID & Negative result; the metric is extra data. \\
Link law for noncompact $G$ (scope paragraph) & Normalizability of a loop-penalty link law for noncompact or general linear $G$. & \OPEN & Declined, with no ledger entry; the program's target case. \\
Three agreements: \MT{prop:geo-three-agreements} & Holonomy characterizations; law coherence implies statistic agreement, not conversely; frame agreement independent; one loop insufficient; $\GL(2,\R)$ witness with non-orthogonal holonomy. & \SOLID & Any $G$; witnesses refute general implications. \\
Agreement and partitions (scope paragraph) & Whether law coherence bears on selecting a partition into meta-agents. & \OPEN & Part of \MT{open:cg-natural-selector}. \\
\end{scorecard}

\begin{hygiene}
Three conditions matter for the models the program cares about. Compactness of $G$ and unitarity of the
representation are what make the loop-penalty link law normalizable; they hold for $SU(N)$ and fail for
$\GL(K,\R)$ and $\GL(K,\mathbb C)$, and the chapter supplies no substitute. The existence of cross morphisms is
additional structure: the bundles do not supply maps between belief laws and model laws, and with irreducible
inequivalent linear representations no nonzero equivariant linear map exists. The single frame shared by the
two channels of an agent is justified only because no energy, kernel, or evaluator reads the relative frame
between them (\MT{prop:geo-common-frame-gauge}); a model that read it would need both frames back. The
common-link specialization $\Theta_e^b=\Theta_e^m$ is an optional modeling choice, not a consequence.

The rest is hygiene here: the context space is a second countable Hausdorff manifold (idle at $c_*$),
representations are measurable, sample fibers are standard Borel, and the smooth tier (\MT{hyp:geo-smooth-tier}),
which excludes infinite-dimensional law fibers, matters only for the fiber geometry of \cref{comp:ch07}.
\end{hygiene}

The chapter has no toy results: its small graphs and its $SO(2)$, $SO(3)$, and $\GL(2,\R)$ instances are
counterexamples to general implications, solid at full generality, or illustrations.

\section{Open ends}
\label{comp:ch02-open}

Two questions are left unsettled. The first is whether a normalizable law on links built from loop penalties
exists for a noncompact or general linear structure group. The chapter declines the claim and records no open
problem or settling condition for it; for the $\GL(K)$ target, links therefore remain structural parameters, and
the theory supplies no Wilson-type prior to integrate them against. The second is whether any of the three agreements bears on
choosing a partition of agents into meta-agents, part of \MT{open:cg-natural-selector}, which is settled by a
relabeling-natural, gauge-compatible functional of the laws, not monotone under refinement, with a nondegenerate
optimizer, or by a theorem that no such selector exists in a stated class. Connections, base transport, and
descent are out of scope by design rather than open.

\begin{verdictbox}{\DELIVERED}
The chapter says that it determines what geometric data an agent carries at $c_*$ and which are invariant: the
active gauge group is one copy of $G$ inside the passive vertex group, every link carries an invariant
conjugacy class with loop holonomy the telescoped product of per-edge twists, and three agreements are
distinct; the introduction's reading map repeats this. Every one of these statements is proved, for an
arbitrary Lie group, with short proofs. The mathematics is standard in substance: torsor algebra, automorphisms
of a torsor, degree-one nonabelian cohomology of a graph, and the holonomy principle on a graph. What is new is
the typing, and it does real work in one place: because all agents share one fiber, links become automorphisms
of that fiber and acquire per-edge invariants that pure lattice gauge theory lacks, visible to a model only
through frame-referenced structure. No connection, curvature, or bundle topology enters, and the text does not
claim otherwise. The gap that matters for the program is the declined normalizability of loop-penalty link laws
for noncompact $G$: for $\GL(K,\R)$ the chapter supplies an invariant penalty only after adding a covariant fiber
metric, and no normalized link law.
\end{verdictbox}


\chapter{Probability Semantics at \texorpdfstring{$c_*$}{c*}}
\label{comp:ch03}

\section{The question and the answer}
\label{comp:ch03-question}

The technical chapter asks on which measurable objects, and in which order, every probabilistic statement at
$c_*$ rests. Its answer is a discipline. Probability laws and Markov kernels are the primitive objects.
Densities are Radon--Nikodym derivatives against reference measures fixed once for a whole family of laws,
generative and recognition alike. Logarithms of density ratios are written only where the numerator is
absolutely continuous with respect to the denominator, and relative entropy is $+\infty$ everywhere else
(\MT{def:prob-log-density-ratio-convention}). The stated main results are that a family of mixed laws on the
real line has a common dominating measure exactly when the union of its atoms is countable
(\MT{thm:prob-mixed-family-common-domination}), the relative-entropy chain rule in $[0,+\infty]$ that every
later decomposition uses (\MT{lem:prob-kl-chain-rule}), and three levels of description, sample, law point,
and configuration, at which a divergence between laws is not an energy of samples
(\MT{prop:prob-law-point-energy}) and a point-mass recognition of a law point has infinite cost
(\MT{prop:prob-law-point-point-mass}). Along the way the chapter types the population recognition law as a
selected coupling of the agents' local laws and prices restricted access. The generative law of
\cref{comp:ch04} is built on these spaces, the evidence bound of \cref{comp:ch05} is an application of the
chain rule, and attention (\cref{comp:ch06}) and formation (\cref{comp:ch13}) are typed by the levels.

Which of this is substance and which is hygiene? Most of the chapter is measure theory done carefully so that
later identities hold without case distinctions, and a physicist can take it on trust. Three things are
substantive. The population recognition law is not determined by the agents' local beliefs and must be
selected, and that selection is where recognition-level correlations between agents live. Restricting what
each agent can see has an evidence price that no selection removes. And the level typing shows that a
relative entropy between two agents' beliefs, the building block of the program's attention free energies, is
not the expectation of any energy of samples, so its place inside one evidence bound depends on a prior over
beliefs that the theory has not constructed (\MT{open:prob-law-point-prior}). A fourth item sits on the border:
families of laws with moving atoms, and noiseless records with a continuum of values, have no density model,
and the evidence bound for such records is open (\MT{open:prob-nondominated-likelihood}).

\section{The picture}
\label{comp:ch03-picture}

\subsection{Laws first, because the theory needs degenerate laws}

A physicist's default is to write a probability as a density and a free energy as an integral of
$q\log(q/p)$. The chapter refuses that default for a concrete reason: the models it must cover include
point-mass model presentations (the definite case $q_i^m=\delta_{m_i}$ of \MT{def:prob-model-evaluation}),
laws carried by lower-dimensional sets such as proper affine subspaces, and laws with atoms whose locations
move with a parameter. None of these has a density against Lebesgue measure, and some families of them have
no common reference at all. A sigma-finite measure has at most countably many atoms, so the moving-atom family
$\mu_t=\tfrac12\mathcal N(0,1)+\tfrac12\delta_t$, $t\in\R$, has no common dominating measure, and neither does
any uncountable family of Dirac laws on an uncountable space (\MT{thm:prob-mixed-family-common-domination},
\MT{sec:exp-dominated-nondominated}). Working with measures and setting relative entropy to $+\infty$ off
absolute continuity keeps every identity meaningful in the extended reals: the evidence bound of a singular
recognition law is $-\infty$ rather than undefined. This is the whole content of "kernels before densities".
It costs nothing when densities exist, and it keeps the degenerate cases inside the theory when they do not.

\begin{analogy}{exact}
Conditioning on an event of probability zero is not determined by the joint law: the posterior at a single
observation depends on the version of the conditional one uses, which is the Borel--Kolmogorov paradox. The
chapter's regular-conditional convention (\MT{hyp:prob-regular-conditional-convention}) therefore asserts
every fixed-observation statement only for almost every observation, and its witness
(\MT{rem:prob-version-witness}) is an instance of the paradox: two versions of one joint density on $\R^2$ give
different posteriors at $o=0$.
\end{analogy}

\subsection{Local beliefs do not fix the population law}

Each agent has a local recognition law $\mathbb Q_{i,o,X}$ on its own block, computed from what it can access.
The population recognition law is a law on all blocks at once whose agent marginals are the local laws, that
is, a member of the coupling class $\operatorname{Cpl}(\mathbb Q_{\bullet,o,X})$. The class is never empty and
is convex (\MT{prop:prob-coupling-existence-convexity}), and it is a single point only when at most one local
law is non-Dirac (\MT{prop:prob-marginals-do-not-determine-joint}). As soon as two agents are uncertain, their
beliefs underdetermine the joint recognition law, and the theory selects one by a declared rule
(\MT{hyp:prob-measurable-coupling-selector}). The product coupling is one admissible selection, not the
default; copulas \citep{Nelsen2006}, shared latent variables, ordered conditionals, or graphical families
supply others. The coupling is a modeling input, not an output of the local beliefs, and the manuscript
asserts no canonical selector.

Restricted access sharpens the point (\MT{prop:prob-access-restriction}). An access map limits what each local
law may depend on, but a selector that reads all records can still make the population law depend on data that
no agent sees, so a claim of decentralized inference also needs a selector indexed by the agents' accesses
alone. And restricting access costs evidence. When an agent cannot distinguish two observations whose
posteriors differ, its single local law sits at a positive divergence from at least one of them; in the
chapter's two-agent binary witness the worst-case gap is at least the Chernoff information between the two
posterior marginals, about $0.0885$ nats \citep[Section~11.9]{CoverThomas2006}.

\begin{analogy}{structural}
The product coupling is the naive mean-field trial state of the Gibbs--Bogoliubov variational principle, and
correlated couplings are richer trial states with the same one-body marginals. The identification is exact for
the product case; for the others it is structural, since a coupling class fixes the one-body marginals and
leaves every higher correlation free.
\end{analogy}

\subsection{Three levels of description}

The chapter distinguishes three kinds of latent variable (\MT{def:prob-description-levels}). At the sample
level the latent state is the population sample $y=((k_i,m_i))_{i\in V}$, the agents' microstates, and an
agent's belief is the marginal on $k_i$ of the recognition law, a law over microstates. At the law-point level
the beliefs themselves are latent: the latent state of agent $i$ is a point $\mathfrak q_i^b$ of the
associated law fiber, a probability law treated as a random variable, and a recognition law is a law over laws.
At the configuration level the latent state is a presentation $w$, for instance a mean and a covariance,
together with a sample $y$ drawn from the law that $w$ presents, the two linked by a presentation kernel. The
levels define different optimization problems, and the manuscript requires every energy to name its level.

The reason is \MT{prop:prob-law-point-energy}. An energy $D$ of samples enters a free energy as $\E_QD$, which
is affine in the recognition law. A transported relative entropy $\KL(q_i\Vert\widehat\rho_x(\Theta_e^x)q)$ is
strictly convex in the source law $q$ along mixtures, because $-\log$ is strictly convex, so no sample energy
can match it at two source laws and one strict mixture of them. A divergence between two agents' beliefs is a
functional of the ensemble, like an entropy, and not a Hamiltonian term. The consequence runs through the
program. The belief-alignment terms $\KL(q_i\Vert(\Theta_e)_\#q_j)$ that weight attention rows
(\MTeq{eq:full-graph-transported-divergences}) are law-point energies, and they become terms of one evidence
bound only if beliefs are latent law points under a normalized prior, which is \MT{open:prob-law-point-prior}.

\MT{prop:prob-law-point-point-mass} closes the obvious shortcut. If law points are latent and their prior has
no atom at the law point one wants, a recognition law that puts all its mass on that law point has infinite
relative entropy to the posterior and an evidence bound of $-\infty$. A row or an update computed at fixed law
points is therefore not a recognition optimum at the law-point level. With law points latent it is a
generative posterior quantity conditional on them; with law points declared as parameters it is a function of
parameters whose change is an M-coordinate subject to the acceptance rule of
\MT{def:elbo-generalized-m-acceptance}. Both readings are legitimate, and neither is an agent's recognition of
its own belief.

\begin{analogy}{structural}
The entropy $-\int q\log q$ is a functional of the ensemble and not the average of any state-independent
microscopic energy: it is strictly concave in $q$, while every average $\int qH$ is linear. The cross-entropy
between two agents' beliefs has the same character in its second argument, and that is all
\MT{prop:prob-law-point-energy} uses. The point-mass result has an exact counterpart: a trial state
concentrated on one point of a continuum has entropy $-\infty$ relative to a continuous reference.
\end{analogy}

\subsection{Auxiliary variables are free only when exact}

The configuration level is reached from the sample level by a lift
(\MT{prop:prob-configuration-lift-identity}). Draw configuration coordinates from a kernel
$\kappa^{\mathrm{cfg}}(y,dw)$ given the sample: the evidence does not change, and a lifted recognition law pays,
relative to the sample-level bound of its sample marginal, the expected divergence between its conditional law
of $w$ given $y$ and $\kappa^{\mathrm{cfg}}$. The lift never raises the bound, and it leaves a finite bound
unchanged exactly when recognition uses $\kappa^{\mathrm{cfg}}$ itself as its conditional.

\begin{analogy}{exact}
Introducing an auxiliary variable with its exact conditional law, as in a Hubbard--Stratonovich or
data-augmentation step, leaves the partition function unchanged, and a variational ansatz on the enlarged
space pays the conditional relative entropy for any mismatch. Both statements are the chain rule of
\MT{lem:prob-kl-chain-rule}.
\end{analogy}

\section{Objects and notation}
\label{comp:ch03-objects}

\begin{decoder}
Sample blocks $\mathsf Y_i=\mathsf K_i\times\mathsf M_i$; latent space $\mathsf Y$ & Agent $i$'s belief-sample and model-presentation coordinates; their product over agents. & Single-site and full configuration space \\
Record $a\in\mathfrak A$; scope $\partial a$; record space $\mathsf O_a$ & An observation factor attached to a set of agents; a single-agent observation has singleton scope. & Measured quantity on a bond or plaquette \\
Structural configuration $X\in\mathsf X$ & The multigraph, links, frames, and mechanism parameters, held fixed by declaration. & Lattice, couplings, background fields \\
Access map $\operatorname{Acc}_i$ & What agent $i$ may use from the records and the structure. & A local observer's data \\
Local recognition kernel $\mathbb Q_i$; law $\mathbb Q_{i,o,X}$ & Agent $i$'s approximate posterior on its own block, computed from its access. & One-body trial state \\
Population recognition law $\mathbb Q_{V,o,X}$ & A joint approximate posterior on all blocks, selected to have the local laws as marginals. & Many-body trial state \\
Coupling class $\operatorname{Cpl}(\mathbb Q_{\bullet,o,X})$ & All joint laws with the given agent marginals. & Fr\'echet class of joint laws \\
Selection rule $\mathsf S$ & A declared measurable rule picking one coupling. & Choice of variational ansatz \\
Recognition marginals $q_i^{b;o,X}$, $q_i^{m;o,X}$ & The belief and model parts of the local law; the values of agent $i$'s sections at $c_*$. & One-body reduced laws \\
Generative law $\mathbb P_{\theta,V}(\cdot\given X)$ & The constructed joint of latents and records (\cref{comp:ch04}). & Prior times likelihood; Boltzmann weight \\
Posterior slice $\boldsymbol\Pi_{\theta,V,o,X}$ & The exact conditional law of latents at observation $o$. & Gibbs state conditioned on data \\
Evidence $p_{\theta,V}(o\given X)$ & The marginal density of the records. & Partition function \\
References $\nu^{\mathsf Y}$, $\nu^{\mathsf O}$, $\nu_{\mathrm{mix}}$ & Fixed measures against which densities are taken; $\nu_{\mathrm{mix}}$ adds countably many atoms to Lebesgue measure. & Phase-space measure \\
Relative entropy $\KL(Q\Vert P)$ & Kullback--Leibler divergence, set to $+\infty$ off absolute continuity. & Free-energy excess over equilibrium \\
Presentation $m_i$; evaluator $\operatorname{ev}_i$ & A model coordinate that the evaluator turns into a generative kernel; not itself a law. & Coupling constants of a local Hamiltonian \\
Law point $\mathfrak q_i^x$; block $\mathsf Y_i^{\mathrm{law}}$ & A law treated as a latent random variable. & A one-body distribution treated as a field \\
Configuration space $\mathsf Y^{\mathrm{cfg}}$; $\operatorname{pres}$; $\kappa^{\mathrm{cfg}}$ & Presentations such as mean and covariance; the map to law points; the kernel linking them to samples. & Auxiliary-field representation \\
Regular set $\mathsf O^{\mathrm{reg}}_{\theta,V,X}$ & Observations of full probability with finite positive evidence, where the posterior formula is used. & None \\
Extended reals & Values in $[-\infty,+\infty]$; bounds may be $-\infty$ and divergences $+\infty$. & None \\
\end{decoder}

Several terms carry nonstandard meaning. "Configuration" is used twice, and the two uses must not be merged: the
structural configuration $X$ collects graph, links, frames, and parameters, while the configuration level of
description has latent presentations $w\in\mathsf Y^{\mathrm{cfg}}$. "Kernel" always means a Markov kernel, a
measurable family of probability laws, never an integral-operator kernel or a positive-definite kernel.
"Record" is an observation factor with a scope, a hyperedge on the agents, and never a memory. "Recognition
law" is the approximate posterior of variational inference, and the exact conditional is called the posterior
slice. "Evidence" is the marginal likelihood, which a physicist should read as a partition function.
"Presentation" is a parameter-like coordinate of a generative factor and has nothing to do with a group
presentation. "Version" means a representative up to a null set. "Selected" means chosen by a declared rule,
not derived.

\section{The formulas, explained}
\label{comp:ch03-formulas}

\paragraph{Relative entropy.}
The convention of \MTeq{eq:prob-kl-definition} is
\[
\KL(Q\Vert P)=
\begin{cases}
\displaystyle\int\log\Bigl(\frac{dQ}{dP}(y)\Bigr)Q(dy), & Q\ll P,\\[1.2ex]
+\infty, & \text{otherwise},
\end{cases}
\]
the divergence of \citet{Kullback1951} with the extension of \citet{Csiszar1967}. The negative part of the
integrand has integral at most $1/e$, so the first branch is a well-defined element of $[0,+\infty]$. The
typical infinite case in the program is a Gaussian with singular covariance measured against one with
nonsingular covariance; the typical finite case is two nondegenerate Gaussians. Whether $Q\ll P$ can be read
off from densities: against a common reference it holds exactly when $Q$ puts no mass where the density of $P$
vanishes (\MT{prop:prob-density-absolute-continuity}).

\paragraph{The mixed reference.}
When the atoms of a family of mixed laws on $\R$ form the countable set $\{a_1,a_2,\ldots\}$,
\[
\nu_{\mathrm{mix}}=\operatorname{Leb}+\sum_{n\geq1}2^{-n}\delta_{a_n}
\]
dominates the whole family with one parameter-independent measure (\MTeq{eq:prob-mixed-family-reference}). If the
atoms move over an uncountable set, no such measure exists, and the chapter then declares no density model on that
coordinate. For learning this is concrete: a law with an atom whose location is a learned parameter cannot have
its densities compared across parameter values against one reference.

\paragraph{The chain rule and its corollaries.}
For laws on a product with first marginals $Q_1,P_1$ and conditional kernels $Q_{2|1},P_{2|1}$,
\[
\KL(Q\Vert P)=\KL(Q_1\Vert P_1)+\int\KL\bigl(Q_{2|1}(y_1,\cdot)\big\Vert P_{2|1}(y_1,\cdot)\bigr)Q_1(dy_1)
\quad\text{in }[0,+\infty]
\]
(\MTeq{eq:elbo-kl-chain-rule}). This is the relative-entropy chain rule, standard in substance
\citep[Section~2.5]{CoverThomas2006}; the chapter's version needs no densities and no finiteness hypothesis, and
it does not depend on the chosen conditional versions. The same lemma gives four further tools
(\MT{lem:prob-kl-chain-rule}). Measured against a product reference, a joint law's divergence splits into its
total correlation $\KL(Q\Vert\bigotimes_bQ_b)$ plus the marginal divergences; physically, the connected
correlations separate from the one-body terms. With the first marginal fixed, the divergence is minimized by
copying the reference's conditional. Divergence cannot increase under a Markov kernel, which is the data
processing inequality (clause (d), also cited as \MT{thm:cg-kl-dpi-extended}). And divergence is unchanged by a
bimeasurable bijection, which is why frame changes preserve every divergence without any density hypothesis.

\paragraph{The evidence price of restricted access.}
In the witness of \MT{prop:prob-access-restriction}, agent $1$ sees only its own record, so its local law
$\operatorname{Ber}(t)$ is the same at the observations $(1,0)$ and $(1,1)$, whose posterior marginals for $y_1$
are $\Pi_1^{(0)}=(37/89,52/89)$ and $\Pi_1^{(1)}=(13/161,148/161)$. Then
\[
\max_{o_2\in\{0,1\}}\KL\bigl(\mathbb Q_{V,(1,o_2),X}\big\Vert\boldsymbol\Pi_{\theta,V,(1,o_2),X}\bigr)
\geq\max_{o_2}\KL\bigl(\operatorname{Ber}(t)\big\Vert\Pi_1^{(o_2)}\bigr)
\geq\sup_{\xi\in[0,1]}\bigl(-\log Z_\xi\bigr)=0.08849\ldots
\]
(\MTeq{eq:prob-access-gap-floor}), with $Z_\xi=\sum_{y_1}\Pi_1^{(0)}(y_1)^\xi\Pi_1^{(1)}(y_1)^{1-\xi}$. The first
step is data processing onto agent $1$'s coordinate, and the second is the Chernoff bound: a single law cannot
be close to two distinct targets at once, and the best compromise is the Chernoff point. The bound is attained.
The equality condition of the evidence bound, recognition equal to the posterior, lies outside the
access-restricted class at one of the two observations.

\paragraph{Transported divergences are not sample energies.}
For a receiver law $q_i$ and two source laws $p_1,p_2$ with finite transported divergences whose transported
densities differ on a set of positive $q_i$-measure,
\[
\KL\bigl(q_i\Vert\widehat\rho_x(\Theta_e^x)(sp_1+(1-s)p_2)\bigr)
<s\KL\bigl(q_i\Vert\widehat\rho_x(\Theta_e^x)p_1\bigr)+(1-s)\KL\bigl(q_i\Vert\widehat\rho_x(\Theta_e^x)p_2\bigr)
\]
for $0<s<1$ (\MT{prop:prob-law-point-energy}). The left side is what the divergence does along a mixture of
source laws, and the right side is what any energy averaged over samples must do. Besides finiteness, the
hypotheses exclude only the case in which the two transported laws have equal densities almost everywhere under
the receiver law.

\paragraph{Point masses on law points.}
At the law-point level, with a posterior dominated by the prior (true when every record kernel has a density
against a fixed reference) and a law point $\bar{\mathfrak q}$ that is not an atom of the prior,
\[
\KL\bigl(\delta_{\bar{\mathfrak q}}\big\Vert\boldsymbol\Pi^{\mathrm{law}}_{\theta,V,o,X}\bigr)=+\infty
\]
(\MTeq{eq:prob-law-point-point-mass}). The singleton $\{\bar{\mathfrak q}\}$ is a posterior-null set carrying all of
the recognition mass. With an atom at $\bar{\mathfrak q}$ the divergence can be finite; the proposition concerns
non-atomic priors, which is the situation for any continuous family of candidate beliefs.

\paragraph{The configuration lift.}
For a lifted recognition law $\mathbb J$ with sample marginal $\bar{\mathbb Q}$,
\[
\Lelbo^{\mathrm{ext,cfg}}_V(\mathbb J;X,o)
=\Lelbo^{\mathrm{ext}}_V(\bar{\mathbb Q};X,o)
-\E_{\bar{\mathbb Q}}\Bigl[\KL\bigl(\mathbb J(\cdot\given Y)\big\Vert\kappa^{\mathrm{cfg}}(Y,\cdot)\bigr)\Bigr]
\]
(\MTeq{eq:prob-configuration-lift-identity}). The first term is the bound of the sample-level problem, the second
is the price of any mismatch between recognition's conditional over configurations and the lift kernel, and both
live in the extended reals.

\section{What is proved, and how solid it is}
\label{comp:ch03-status}

\begin{scorecard}
Kernels first; levels: \MT{def:prob-normalized-kernels}, \MT{def:prob-description-levels} & Laws and kernels are primitive, densities are derivatives against family-wide references, and every energy names its level: sample, law point, or configuration. & \CHOICE & A typing discipline; it costs nothing when densities exist. \\
Common domination: \MT{thm:prob-mixed-family-common-domination} & A family of mixed laws on $\R$ without singular-continuous parts has a sigma-finite common dominating measure exactly when its atoms are countable; moving atoms are not dominated. & \SOLID & Standard in substance (a sigma-finite measure has countably many atoms). \\
Dirac families: \MT{sec:exp-dominated-nondominated} & An uncountable family of Dirac laws on an uncountable standard Borel space has no common sigma-finite dominating measure. & \SOLID & Standard. \\
Chain rule and corollaries: \MT{lem:prob-kl-chain-rule} & Chain rule, product-reference split, fixed-marginal minimization, data processing, and bijection invariance, all in $[0,+\infty]$. & \SOLID & Standard in substance; proved in \MT{app:measure-lemmas} without density or finiteness hypotheses. \\
Couplings exist: \MT{prop:prob-coupling-existence-convexity} & The coupling class is nonempty and convex; the product coupling is a measurable kernel. & \SOLID & Standard. \\
Marginals do not fix the joint: \MT{prop:prob-marginals-do-not-determine-joint} & The coupling class is a singleton exactly when at most one marginal is non-Dirac. & \SOLID & Standard in substance (Fr\'echet classes). \\
Coupling selection: \MT{hyp:prob-measurable-coupling-selector} & The population recognition law is chosen by a declared measurable rule inside the coupling class. & \CHOICE & No canonical selector; equivariance of the selector under frame changes is declined. \\
Access restriction, part (i): \MT{prop:prob-access-restriction} & Access maps constrain local laws, not the population law; decentralization also needs an access-indexed selector. & \SOLID & A general no-go shown by counterexample. \\
Access restriction, part (ii): \MT{prop:prob-access-restriction} & Every access-restricted recognition law has worst-case gap at least $0.08849\ldots$ nats, and the bound is attained. & \TOY & Two binary agents with one prior; the constant is instance specific. \\
Law-point energy: \MT{prop:prob-law-point-energy} & A transported relative entropy is strictly convex in the source along mixtures, so no sample energy reproduces it. & \SOLID & Needs finite divergences and transported laws that differ on the receiver's support. \\
Point-mass cost: \MT{prop:prob-law-point-point-mass} & Recognition by a point mass at a non-atomic law point has infinite divergence and bound $-\infty$; rows at fixed law points are posterior quantities or parameter functions. & \SOLID & Needs the posterior dominated by the prior, true for record kernels with densities. \\
Configuration lift: \MT{prop:prob-configuration-lift-identity} & Lifting leaves the evidence unchanged and lowers the bound by the expected conditional divergence. & \SOLID & Standard in substance (augmentation and chain rule). \\
Law-point prior: \MT{open:prob-law-point-prior} & A normalized generative law over law points for a general class. & \OPEN & Needed for transported divergences to be evidence-bound terms. \\
Nondominated likelihoods: \MT{open:prob-nondominated-likelihood} & Evidence bound and coarse channels for record families with no common dominating measure. & \OPEN & Covers noiseless records with a continuum of values. \\
\end{scorecard}

\begin{hygiene}
Four conditions can fail in the models the program cares about. Common domination fails for families with
moving atoms and for deterministic records with a continuum of values, and then the density formulas, and
everything that uses them, are not available (\MT{open:prob-nondominated-likelihood}). Absolute continuity of
recognition with respect to the posterior is not assumed; its failure is the singular branch on which the bound
is $-\infty$, and it genuinely occurs for degenerate Gaussian beliefs and point masses. The coupling selection is
a modeling commitment whose choice changes every collective quantity, and its equivariance under frame changes
is not claimed. The access structure, when imposed, carries an evidence price.

The rest is hygiene a physicist may skip on first reading: standard Borel spaces, sigma-finite references, joint
measurability of Radon--Nikodym derivatives and of relative entropies in their parameters (proved once in
\MT{app:measure-lemmas}), the almost-every-observation convention for posteriors, and bookkeeping in the extended
reals. The restriction of the common-domination theorem to laws without singular-continuous parts is also
hygiene: no model in the program uses such laws.
\end{hygiene}

\section{Toy versus general}
\label{comp:ch03-toy}

The only toy result is the quantitative half of the access-restriction proposition. Its mechanism is visible in
the proof and has a general form: data processing onto the coordinates an agent can see, followed by the Chernoff
bound for one law facing two targets. Nothing in that argument depends on binary variables or on two agents, so
the mechanism may well survive in large networks with many agents sharing an access class, but the manuscript
states and proves it only for the displayed instance, and the constant $0.08849\ldots$ and the optimizers
$\xi\approx0.548$ and $t\approx0.783$ are specific to that prior and those record accuracies. The qualitative
half, that access restriction constrains local laws and not the population law, is a general negative statement
proved by counterexample.

\section{Open ends}
\label{comp:ch03-open}

Two open problems are recorded. The first is a normalized prior over law points (\MT{open:prob-law-point-prior}).
Law fibers carry no canonical sigma-finite reference, since the space of all laws on a sample fiber is infinite
dimensional, and the Haar measure of a noncompact $G$ is infinite, so a Haar prior on a free orbit does not
normalize. The problem is settled, for a class of law points, by a sigma-finite reference on the law-point or
presentation space, a finite normalizer for the exponentiated law-point energies, the transformation law of the
prior under the active gauge action, record kernels that read law points, and a compatibility relation between
the law-point posterior and the sample-level recognition marginals. Until then, the transported divergences
\MTeq{eq:full-graph-transported-divergences} enter as quantities conditional on law points, as functions of
parameters, or at the configuration level under the one proper prior the formation model supplies. The second is the treatment of record families with no
common dominating measure (\MT{open:prob-nondominated-likelihood}): the joint law and its posterior exist, but the
likelihood and possibly the evidence density do not. It is settled by a treatment of the evidence bound, its
local disintegrations, and the coarse channels at the level of joint measures and regular conditional laws in
which the exact identities of the later chapters are proved, or by a record family for which they provably have
no measure-level counterpart.

The scope paragraph also records three limits. Injectivity of a model evaluator is not assumed, and no regularity
of the space of presentations modulo equal evaluated kernels is supplied, so presentations are kept as they are;
this is bookkeeping for overparameterized models. The selector's equivariance under frame re-expression is not
claimed; this is substantive, because a recognition procedure rerun in new frames need not return the
transformed coupling, and the invariance results of \cref{comp:ch04} are stated for a coupling transported as a
whole.

\begin{verdictbox}{\DELIVERED}
The chapter says that laws and kernels come first and densities follow against family-wide references, so that
the evidence bound becomes an identity in the extended reals, and names as its main results common domination
exactly when atoms are countable, the chain rule in $[0,+\infty]$, and three levels at which a law-point energy
is not a sample energy and point-mass recognition of law points has infinite cost; the reading map adds that it
proves the chain rule used everywhere and distinguishes the three levels. All of this is proved at the generality
the theory uses. Most of it is standard measure theory and information theory in careful form, and a physicist
loses nothing by treating it as such. The substantive content is elsewhere and is also delivered: local beliefs
do not determine the population recognition law, so a coupling must be selected; restricted access has a
provable evidence price; and a divergence between beliefs is not a sample energy. The last result exposes a gap
that the chapter names rather than closes: without a normalized law-point prior, the belief-alignment divergences
at the center of the program's attention free energies are not terms of one evidence bound at the law-point
level.
\end{verdictbox}


\chapter[The Fixed Normalized Generative Law]{The Fixed Normalized\\Generative Law}
\label{comp:ch04}

\section{The question and the answer}
\label{comp:ch04-question}

The technical chapter asks how one normalized joint law of latent samples and records is built at $c_*$, and
what a passive re-expression of the frames does to it. The construction is directed. Structural data fix a
directed acyclic graph $\Gamma$ on the agents. Each agent has a model factor, the law of its model presentation
given its parents' models, and a state factor, the law of its belief sample given its parents' beliefs and its
own model; composed, they form the agent's mechanism, and composing the mechanisms along a topological order of
$\Gamma$ gives the latent law (\MT{def:gen-parameter-split}, \MT{cons:gen-finite-directed-law}). Records are then
attached as conditionally independent children of the latents, each with an arbitrary scope. The chapter proves
that the result is normalized for every finite record hypergraph, cyclic or not (\MT{prop:gen-exact-normalization});
that locally normalized reciprocal couplings between latents can fail to give any joint
(\MT{prop:gen-gibbs-counterexample}); and that, under a covariance hypothesis on the kernels, a passive
re-expression pushes the whole joint forward, so the evidence and the bound of a transported recognition law are
unchanged (\MT{prop:gen-product-evidence-invariance}, \MT{cor:elbo-independent-frame-invariance}). It also states,
as a standing hypothesis, the typing prohibition: no generative factor may read a recognition law or the
posterior (\MT{req:gen-typing-prohibition}). The evidence of this joint is bounded in \cref{comp:ch05}, and
attention and coarse channels act on it in \cref{comp:ch06} and \cref{comp:ch09}.

What is nontrivial about normalizing on graphs with cycles? Mathematically, nothing: the proof integrates out the
records first, each has unit mass, and what remains is a directed composition of normalized kernels. This is the
standard fact that a Bayesian network is normalized whatever the undirected loops of its moral graph
\citep{Lauritzen1996}. The content lies in what the construction requires. Generative influence among latents
must be acyclic, and every interaction that closes a loop must be written as a record, a normalized law of an
observed quantity given the latents it reads. The alternative, reciprocal couplings written directly between
latents, has a parameter-dependent normalizer that can be infinite even when every local piece is normalized,
which is what the Gaussian counterexample shows. The chapter chooses the directed route because its
normalization does not depend on the parameters, and the price is that loop interactions enter only through
conditioning on records.

\section{The picture}
\label{comp:ch04-picture}

\subsection{Three graphs, not one}

A physicist will be tempted to see one interaction graph. There are three, with different roles. The directed
acyclic graph $\Gamma$ says who generates whom among latents. The interaction multigraph $\mathfrak I$ of
\cref{comp:ch02} carries the links. The record hypergraph, formed by the scopes $\partial a$, says which latents
each observation reads. The running example makes the distinction concrete: $\Gamma$ is a tree in which agent $1$
is the parent of agents $2$ and $3$, $\mathfrak I$ is the triangle, and two records with scopes $\{1,2\}$ and
$\{2,3\}$ close the latent tree into the triangle. Cycles live in the record hypergraph; $\Gamma$ never has any.

\subsection{Records close loops, and the evidence is the partition function}

When each record kernel has the form $K_a(X,y_{\partial a},do_a)=\exp(-E_{a,o}(y_{\partial a}))\nu_a(do_a)$, the
joint density is the latent density times $\exp(-\sum_aE_{a,o})$ (\MTeq{eq:gen-record-density}). Combined with
the posterior formula of \MT{hyp:prob-regular-conditional-convention}, this says that the posterior is a Gibbs
measure whose reference is the directed prior and whose Hamiltonian is the sum of the record energies, and that
its partition function is the evidence. A binary record whose success probability is $\exp(-E)$ with $E\geq0$,
observed at success, contributes exactly the Boltzmann factor $\exp(-E)$ to the posterior, and the running example
uses this device for its loop-closing records. Normalization of the joint is automatic; what can fail is
finiteness and positivity of the evidence, and the manuscript admits only observations with finite positive
evidence (\MT{sec:intro-conventions}). One physical consequence is easy to miss. Sampling the latents from the
generative model ignores the records altogether, because the records are children of the latents, not
constraints on them; the loop interaction shapes the latents only after conditioning on the observed records.

\begin{analogy}{exact}
Prior times likelihood, normalized by the evidence, is a Gibbs measure with the prior as reference and the
record energies as Hamiltonian; the evidence is its partition function. This is Bayes' rule and needs no further
argument.
\end{analogy}

\subsection{Why the loop is not a potential between latents}

The undirected alternative writes the joint as a product of potentials on agents, edges, and records, divided by
a partition function $Z_X$, and it is legitimate whenever $0<Z_X<\infty$ (\MT{def:gen-undirected-model}). The
trouble is that $Z_X$ depends on the parameters and nothing local controls it. Two scalar agents with normalized
Gaussian site potentials and the positive, finite, continuous edge potential $e^{\mathrm ck_1k_2}$ give the
quadratic form with matrix $\mathsf W_{\mathrm c}$, whose smaller eigenvalue $1-\mathrm c$ is not positive for
$\mathrm c\geq1$, and the integral diverges (\MT{prop:gen-gibbs-counterexample}). Each conditional
$k_1\given k_2$ and $k_2\given k_1$ is a perfectly normalized Gaussian, and still there is no joint.

\begin{analogy}{exact}
This is the Gaussian spin model beyond its critical coupling: once the ferromagnetic coupling exceeds the
single-site stiffness, the quadratic form loses positivity along the uniform mode and the partition function
diverges. The counterexample is the two-site case of the same integral.
\end{analogy}

A related temptation is to append an energy $U$ to the free energy of a joint that does not contain it. The
chapter's remark computes the result (\MT{rem:gen-appended-energy}): for $U$ bounded below,
$\Fenergy_V[Q;X,o]+\E_QU$ equals the free energy of the tilted joint $Z_U^{-1}e^{-U}\mathbb P_{\theta,V}$ minus
$\log Z_U$ (\MTeq{eq:gen-appended-energy}). Because $\log Z_U$ does not depend on $Q$, recognition updates under
the appended objective are those of the tilted model; because it depends on $(\theta,X)$, comparisons across
parameter values, and hence learning, are off by $\log Z_U$ unless it is restored.

\subsection{The typing prohibition}

The standing hypothesis \MT{req:gen-typing-prohibition} states that every generative factor and record kernel is a
function only of the parameters $\theta$, the structural data $X$, and its displayed latent coordinates. No
generative kernel takes as an argument the recognition law, one of its marginals or parameters, or the posterior.
In the two-argument form $\Fenergy_V(\mathbb Q;\theta)$ of the free energy \citep{Neal1998}, every evaluation holds
the generative argument fixed while the recognition argument varies, which excludes the composite
$\mathbb Q\mapsto\Fenergy_V(\mathbb Q;\theta(\mathbb Q))$ inside one evaluation. The chapter's example of what is
forbidden is a transition built from a live peer recognition marginal, such as agent $i$'s prior set to the
current transported belief of agent $j$. What is allowed is coupling through latents (agent $i$'s state factor
reads agent $j$'s sample), a parameter that depends on the observation or on an earlier recognition law provided
it is held fixed while the bound is evaluated, slow evolution of an agent's model as recognition dynamics of the
model channel or as accepted M-coordinates, each accepted value defining a new fixed joint, and loops of inference
updates under a fixed joint.

Why forbid it? The manuscript's answer is \MT{prop:gen-no-distinguished-target}: a family of joints indexed by
recognition laws satisfies the evidence identity for each member separately, but the change of the bound between
two members mixes a change of evidence with a change of gap, and the family designates no joint as the model. A
companion witness shows that the stronger failures one might expect, such as evidence that must vary along the
family, need not occur (\MT{prop:gen-moving-target-witness}). The prohibition is therefore a commitment to having
one target, not the avoidance of a contradiction.

\begin{analogy}{structural}
In naive mean-field theory for the Ising model, the Gibbs--Bogoliubov functional $\E_QH-TS[Q]$ over product trial
laws uses the fixed Hamiltonian $H=-J\sum_{\langle ij\rangle}s_is_j$, which reads samples. Scoring $Q$ against the
effective Hamiltonian $H_{\mathrm{eff}}(s;m)=-J\sum_i\bigl(\sum_{j\sim i}m_j\bigr)s_i$ with the current magnetizations
$m=\E_Qs$ inserted counts every bond twice, which is why the mean-field free energy carries its familiar correction
$+J\sum_{\langle ij\rangle}m_im_j$. Holding $m$ fixed during each variation and updating it between variations gives
the same stationarity equations. The first procedure is the composite the typing prohibition excludes; the second is
the admissible parameter held fixed while the bound is evaluated. The illustration is the companion's, not the
manuscript's.
\end{analogy}

\subsection{Passive re-expression moves the whole joint}

A passive re-expression acts on latent samples by the block map $R=\bigoplus_i(R_i^b\oplus R_i^m)$, leaves the
records alone, and transforms the frames and links in $X$ by the vertex law. The kernel-covariance hypothesis
(\MT{hyp:gen-kernel-covariance}) asks that re-expressed parameters $(\theta',X')$ exist for which every factor is
the pushforward of the old one and every record kernel reads the rotated latents as it read the old ones. Under
it, the new joint is the pushforward of the old one as a measure, with no density hypothesis, the observation
marginals agree, and the evidence is unchanged for almost every observation
(\MT{prop:gen-product-evidence-invariance}). Relative entropy is invariant under bijections, so a recognition law
transported as a whole keeps its divergence to the transported posterior and its extended bound, singular branch
included (\MT{cor:elbo-independent-frame-invariance}).

\begin{analogy}{exact}
This is a change of integration variables in a partition function with gauge-invariant observables: the records
are the observables, and the Jacobian $\lvert\det R\rvert^{-1}$ appears in densities but cancels from every ratio,
which is why the measure-level statement needs no density hypothesis.
\end{analogy}

\section{Objects and notation}
\label{comp:ch04-objects}

\begin{decoder}
Structural data $X$; parameters $\theta$ & $X$ holds the directed graph, the multigraph with links, and the frames; $\theta$ holds the numerical parameters that learning may move; the split is declared. & Lattice and background fields; couplings \\
Directed acyclic graph $\Gamma$; parents $\operatorname{pa}(i)$ & Who generates whom among latents. & A causal ordering; none in equilibrium statistical mechanics \\
Model factor $P^m_{\theta,i}$ & Law of agent $i$'s model presentation given its parents' models. & Hierarchical prior on couplings \\
State factor $P^k_{\theta,i}$; evaluated kernel $K^X_{i,m_i}$ & Law of agent $i$'s belief sample given its parents' beliefs and its own model; the same kernel indexed by the model sample. & Conditional Boltzmann weight at fixed couplings \\
Agent mechanism $G_{\theta,i}$ & Model factor then state factor, one normalized kernel from the parents' block to the agent's block. & Local stochastic update rule \\
Root; law bridge $\widetilde\Phi^{\mathrm{gen}}_i$ & An agent without parents; the belief law induced by a law over its model. & Posterior-predictive law \\
Directed latent law $\mathbb P^{\mathsf Y}_{\theta,V}$ & Ordered composition of the mechanisms along a topological order. & Bayesian-network prior \\
Record kernel $K_a$; energy $E_{a,o}$ & Normalized law of record $a$ given the latents in its scope; its negative log density at the observed value. & Likelihood factor; bond energy of the posterior \\
Joint $\mathbb P_{\theta,V}(dy,do\given X)$ & The directed latent law times the record kernels. & Prior times likelihood \\
Undirected model; potentials $\psi$; $Z_X$ & A Gibbs model on agents, edges, and records with its partition function. & Boltzmann distribution \\
Appended energy $U$; tilted joint; $Z_U$ & An energy added to a free energy without being part of the joint; the joint it would define; its normalizer. & A potential added without renormalizing \\
Typing prohibition & No generative kernel reads a recognition law, its marginals or parameters, or the posterior. & Hamiltonian independent of the trial state \\
Two-argument free energy $\Fenergy_V(\mathbb Q;\theta)$ & Free energy as a function of recognition law and parameters, varied one at a time. & Gibbs--Bogoliubov functional \\
M-coordinate; accepted value & A change of generative parameters accepted between evaluations. & M-step of expectation maximization \\
Maps $R$, $T$; gauge inert & The passive action on latents, and on latents and records; records are unchanged. & Change of variables; gauge-invariant observables \\
Kernel covariance & The model's coordinate form is compatible with re-expressed parameters. & Gauge covariance of the action \\
Reference factor $j_R$ & Radon--Nikodym factor of the reference under $R$; $\lvert\det R\rvert^{-1}$ for invertible linear blocks. & Jacobian \\
Transported coupling $R_\#\mathbb Q_{V,o,X}$ & A recognition law moved as a whole to the new frames. & Gauge-transformed trial state \\
\end{decoder}

The word "fixed" in the chapter title means held fixed while the evidence bound is evaluated; the joint changes
between evaluations when an M-coordinate is accepted, and each accepted value defines a new joint. "Model" has two
meanings: the agent's model channel ($m_i$ and its law $q_i^m$), and the generative model $\mathbb P_{\theta,V}$.
The model and state factors are generative factors and are not components of the agent of \cref{comp:ch02}, even
though they carry the agent's index; the agent's belief $q_i^b$ is a recognition object. "Cycle" refers to the
record hypergraph; $\Gamma$ is acyclic by definition. "Covariance" in the hypothesis name means gauge covariance,
not a second moment. The evidence invariance holds for almost every observation, not pointwise.

\section{The formulas, explained}
\label{comp:ch04-formulas}

\paragraph{The agent mechanism.}
\[
G_{\theta,i}\bigl(d(k_i,m_i)\given y_{\operatorname{pa}(i)},X\bigr)
=P_{\theta,i}^{m}(dm_i\given m_{\operatorname{pa}(i)},X)K^X_{i,m_i}(k_{\operatorname{pa}(i)},dk_i)
\]
(\MTeq{eq:gen-agent-mechanism}). The model presentation $m_i$ is drawn first from its parents' models, and the
belief sample $k_i$ then from its parents' beliefs using the kernel that $m_i$ evaluates to. At a root the first
factor is a proper law and the second a bridge from the agent's own model to its belief; averaging the bridge over
a model law $s$ gives the belief law $\widetilde\Phi^{\mathrm{gen}}_i(s)$ (\MTeq{eq:gen-root-law-bridge}), which may
be declared to represent the cross morphism of \cref{comp:ch02}. At a nonroot the evaluated kernel is indexed by
the parent beliefs and is not a map of law fibers.

\paragraph{The joint and its normalization.}
\[
\mathbb P_{\theta,V}(dy,do\given X)
=\mathbb P_{\theta,V}^{\mathsf Y}(dy\given X)\prod_{a\in\mathfrak A}K_a(X,y_{\partial a},do_a),
\qquad
\mathbb P_{\theta,V}(\mathsf Y\times\mathsf O\given X)=1
\]
(\MTeq{eq:gen-record-extension}, \MTeq{eq:gen-normalization}). The latent law is the ordered composition of the
mechanisms along a topological order (\MTeq{eq:gen-composition}), and each record kernel has unit mass in its own
coordinate, so integrating the records first removes the product whatever the overlaps of the scopes. Each record
is owned once by the joint, while the local free energies of \cref{comp:ch05} include every factor incident to a
coordinate, as a bond energy is counted once in the Hamiltonian but enters the local field of both of its sites.

\paragraph{The posterior as a Gibbs measure.}
Rewriting \MTeq{eq:gen-record-density} with the posterior formula of \MT{hyp:prob-regular-conditional-convention},
on the regular set of observations,
\begin{gather*}
\boldsymbol\Pi_{\theta,V,o,X}(dy)
=\frac{\exp\bigl(-\sum_{a}E_{a,o}(y_{\partial a})\bigr)\mathbb P^{\mathsf Y}_{\theta,V}(dy\given X)}{p_{\theta,V}(o\given X)},\\
p_{\theta,V}(o\given X)=\int\exp\Bigl(-\sum_{a\in\mathfrak A}E_{a,o}(y_{\partial a})\Bigr)\mathbb P^{\mathsf Y}_{\theta,V}(dy\given X).
\end{gather*}
With no records the posterior is the prior and the evidence is $1$. Cycles in the scopes make the posterior's
factor graph loopy, which matters for computing it and not at all for its normalization.

\paragraph{The running example.}
Agent $1$ is the root, with $k_1\given m_1\sim\mathcal N(0,e^{-m_1}I_3)$; agents $2$ and $3$ are its children, with
$k_2\given k_1,m_2\sim\mathcal N(\rho_b(\Theta^b_{\bar e_1})k_1,e^{-m_2}I_3)$ and
$k_3\given k_1,m_3\sim\mathcal N(\rho_b(\Theta^b_{e_3})k_1,e^{-m_3}I_3)$, and independent standard normal model
coordinates. The record on $\{2,3\}$ has success probability $\exp(-\tfrac12\lvert k_2-\rho_b(\Theta^b_{e_2})k_3\rvert^2)$.
At the conditional means of the children its residual is
\[
\rho_b(\Theta^b_{\bar e_1})k_1-\rho_b(\Theta^b_{e_2}\Theta^b_{e_3})k_1
=\rho_b(\Theta^b_{e_1})^{-1}\bigl(I_3-\rho_b(\mathrm{Hol}^b(p))\bigr)k_1,
\]
whose length is $\sqrt3$ times the component of $k_1$ orthogonal to the rotation axis of the holonomy (the
running-example paragraph after \MT{prop:gen-exact-normalization}). A record that closes the loop measures how far
going around the loop rotates the root's state; the records see the holonomy through an invariant, and with
trivial holonomy the residual vanishes.

\paragraph{The Gaussian counterexample.}
\begin{gather*}
\psi_i(k_i,m_i;X)=(2\pi)^{-1/2}e^{-k_i^2/2},
\qquad
\psi_{12}(y_1,y_2;X)=e^{\mathrm ck_1k_2},
\qquad
\mathrm c\geq1,\\
\mathsf W_{\mathrm c}=\begin{psmallmatrix}1&-\mathrm c\\-\mathrm c&1\end{psmallmatrix},
\qquad
\operatorname{spec}\mathsf W_{\mathrm c}=\{1-\mathrm c,1+\mathrm c\}
\end{gather*}
(\MTeq{eq:gen-gibbs-counterexample}, \MTeq{eq:gen-gibbs-counterexample-matrix}). The product of potentials is
$\exp(-\tfrac12k^\top\mathsf W_{\mathrm c}k)$ up to a constant; in the sum and difference coordinates the integral
factorizes, and the sum direction, with eigenvalue $1-\mathrm c\leq0$, diverges. For $0\leq\mathrm c<1$ the same
model is a legitimate joint.

\paragraph{The appended energy.}
\[
\Fenergy_V[Q;X,o]+\E_QU=\widetilde\Fenergy_V[Q;X,o]-\log Z_U,
\qquad
Z_U=\int e^{-U(y)}\mathbb P_{\theta,V}(dy,do\given X)\in(0,\infty)
\]
(\MTeq{eq:gen-appended-energy}). The left side is what one computes when an energy is simply added; the right side
is the free energy of the tilted joint, which does contain $U$, corrected by its normalizer. The identity holds on
the classical domain with $\E_Q\lvert U\rvert<\infty$.

\paragraph{Covariance.}
Under kernel covariance,
\begin{gather*}
\mathbb P_{\theta',V}(\cdot\given X')=T_\#\mathbb P_{\theta,V}(\cdot\given X),
\qquad
p_{\theta',V}(o\given X')=p_{\theta,V}(o\given X),\\
p_{\theta',V}(o,y'\given X')=\lvert\det R\rvert^{-1}p_{\theta,V}(o,R^{-1}y'\given X),
\end{gather*}
the second for almost every $o$ and the third for invertible linear blocks with Lebesgue references
(\MTeq{eq:gen-gauge-joint-pushforward}, \MTeq{eq:gen-gauge-evidence-invariance},
\MTeq{eq:gen-gauge-lebesgue-density}). For orthogonal representations $\lvert\det R\rvert=1$ and densities simply
move; for $\GL(K)$ they pick up the Jacobian, while measures, evidence, and divergences do not. The belief and model
channels may have different dimensions, and the reference must be the same within each channel. The root law
bridges transform as cross-morphism representatives do, which ties the generative construction to the geometry of
\cref{comp:ch02}. The corollary then gives
\[
\KL(\mathbb Q'_{V,o,X'}\Vert\boldsymbol\Pi_{\theta',V,o,X'})=\KL(\mathbb Q_{V,o,X}\Vert\boldsymbol\Pi_{\theta,V,o,X}),
\qquad
\Lelbo'^{\mathrm{ext}}_V(\mathbb Q'_{V,o,X'};X',o)=\Lelbo^{\mathrm{ext}}_V(\mathbb Q_{V,o,X};X,o)
\]
for $\mathbb Q'_{V,o,X'}=R_\#\mathbb Q_{V,o,X}$ (\MTeq{eq:elbo-gauge-invariance}), the recognition law being
selected before the re-expression and transported as a whole.

\section{What is proved, and how solid it is}
\label{comp:ch04-status}

\begin{scorecard}
Directed construction: \MT{def:gen-parameter-split}, \MT{cons:gen-finite-directed-law} & Latents are generated along a declared acyclic graph, model before belief within each agent; the split between structural data and parameters is declared. & \CHOICE & Excludes reciprocal latent influence; loops must be records. \\
Normalization: \MT{prop:gen-exact-normalization} & Latent composition is a probability kernel; with records conditionally independent given the latents, the joint is normalized for every finite record hypergraph, and its density is the latent density times $\exp(-\sum_aE_{a,o})$. & \SOLID & Standard in substance (Bayesian networks with observed children). \\
Local normalization fails: \MT{prop:gen-gibbs-counterexample} & Normalized agent potentials and a positive, finite, continuous edge potential need not give a finite partition function. & \SOLID & A general negative statement; the two-site Gaussian witness is textbook. \\
Undirected alternative: \MT{def:gen-undirected-model} & A Gibbs model with its partition function is a legitimate joint whenever $0<Z_X<\infty$. & \CHOICE & Admissible, but the finiteness check is parameter dependent and no general criterion is given. \\
Appended energy: \MT{rem:gen-appended-energy} & Adding a bounded-below energy to a free energy equals the free energy of the tilted joint minus $\log Z_U$. & \SOLID & Elementary; the consequence is for learning, since $\log Z_U$ depends on $(\theta,X)$. \\
Typing prohibition: \MT{req:gen-typing-prohibition} & No generative factor or record kernel reads a recognition law, its marginals or parameters, or the posterior; the composite $\Fenergy_V(\mathbb Q;\theta(\mathbb Q))$ is excluded inside one evaluation. & \CHOICE & A standing hypothesis; it admits fixed parameters updated between evaluations. \\
No distinguished target: \MT{prop:gen-no-distinguished-target}, \MT{prop:gen-moving-target-witness} & Recognition-indexed joints obey the identity member by member but designate no target; the stronger failures need not occur. & \SOLID & Justifies the prohibition as a choice of one target, not as avoidance of inconsistency. \\
Kernel covariance: \MT{hyp:gen-kernel-covariance} & Re-expressed parameters exist for which every factor and record kernel transforms covariantly. & \CHOICE & Excludes mechanisms compatible with no re-expression of their parameters. \\
Joint covariance: \MT{prop:gen-product-evidence-invariance} & The re-expressed joint is the pushforward of the old one; evidence is invariant for almost every observation; densities pick up $j_R$, equal to $\lvert\det R\rvert^{-1}$ for linear blocks; root bridges transform as cross-morphism representatives. & \SOLID & Standard in substance (change of variables). \\
Bound invariance: \MT{cor:elbo-independent-frame-invariance} & For a coupling transported as a whole, the posterior transforms by pushforward and the divergence and extended bound are unchanged. & \SOLID & A coupling reselected by a rule in the new frames is not covered; selector equivariance is declined. \\
\end{scorecard}

\begin{hygiene}
Conditions that matter: the latent graph $\Gamma$ must be acyclic, which excludes reciprocal generative influence
and forces loop interactions into records. Records must be conditionally independent given the latents; correlated
record noise has to be modeled by an added latent coordinate. Each record kernel must be a normalized law of its
record for every latent value, which is what makes normalization automatic and what an arbitrary pairwise
potential is not. The typing prohibition and kernel covariance are modeling commitments, the first excluding
peer-belief priors, the second excluding mechanisms whose coordinate form fixes a frame. Finite positive evidence
at an admitted observation is a standing default; it is automatic for binary records whose success probabilities
are positive on a set of positive prior mass, and must be checked for density records. For the undirected
alternative, finiteness of $Z_X$ is a parameter-dependent condition that can fail.

Hygiene: measurability of the kernels and of their composition, Tonelli's theorem, jointly measurable log
densities for the density formula, and the reference condition $R_\#\nu^{\mathsf Y}\ll\nu^{\mathsf Y}$ for the
density form of covariance, automatic for invertible linear blocks with Lebesgue reference.
\end{hygiene}

The chapter has no toy results: the two-site Gaussian counterexample refutes a general implication, and the
running example is an illustration.

\section{Open ends}
\label{comp:ch04-open}

The chapter records no open problem of its own, and its scope paragraph declines two things. Cross morphisms
supply no generative factor, and a model with factors in both directions between the belief and model channels
needs its own normalized joint, which the chapter does not construct; the directed order, model before belief, is
part of the construction rather than a convenience. No quotient of the parameter space by passive re-expression is
formed, and no parameter count subtracts the dimension of the vertex group, because its action can have
nontrivial stabilizers. A physicist used to counting physical degrees of freedom as total minus gauge should note
that the theory does not make this count. Two further limits follow from the construction without being flagged
as open. Reciprocal couplings between latents are admitted only through the undirected model with a verified
finite partition function, and no general criterion for that finiteness is supplied. And the one way around the
typing prohibition that keeps a single target, promoting beliefs to latent law points so that a generative factor
may read them, requires the normalized law-point prior that is \MT{open:prob-law-point-prior}.

\begin{verdictbox}{\DELIVERED}
The chapter says that it builds one normalized joint law of latent samples and records at $c_*$, normalized for
every finite record hypergraph, cyclic ones included; that local normalization of reciprocal latent conditionals
supplies no joint; and that under kernel covariance the evidence and the bound of a transported coupling are
invariant. The reading map adds that it states the typing prohibition. All of this is proved or stated as written.
The normalization and covariance results are standard in substance: a Bayesian network with observed children is
normalized whatever the loops of its moral graph, and a change of variables with gauge-invariant observables
preserves evidence and relative entropy. The substance is in the commitments the chapter makes explicit:
generative influence is acyclic, loop interactions enter as records and therefore act only after conditioning,
reciprocal latent couplings need a separate finiteness check that can fail, and no factor may read a recognition
law. The last commitment is justified as a choice of one target rather than by an inconsistency in the
alternative, which is what the cited propositions actually show. Two narrowings should be kept in view: "cycles"
means cycles of records, not of generative influence, and the invariance of the bound is for a recognition law
transported as a whole, not one recomputed by a selection rule in new frames.
\end{verdictbox}
